\chapter{مقدمه و انگیزهٔ طرح}
%=====================================================================

\section{زمینه و اهمیت موضوع}

امنیت زیرساخت‌های ارتباطی و مالی جهان امروز تقریباً به‌طور کامل بر رمزنگاری کلید عمومی مبتنی بر مسائل ریاضی سخت -- نظیر تجزیهٔ اعداد صحیح بزرگ \lr{(RSA)} و لگاریتم گسسته روی منحنی بیضوی \lr{(ECC/ECDH)} -- استوار است. امنیت این الگوریتم‌ها نه بر یک اثبات ریاضی قطعی، بلکه بر این فرض عملی متکی است که هیچ الگوریتم کلاسیک کارآمدی برای حل این مسائل در زمان چندجمله‌ای وجود ندارد. الگوریتم شور\LTRfootnote{Peter W. Shor} \lr{\cite{shor1997}} نشان داد که یک رایانهٔ کوانتومی با مقیاس کافی می‌تواند این مسائل را در زمان چندجمله‌ای حل کند و بدین ترتیب کل زیرساخت رمزنگاری کلید عمومی کنونی را بی‌اثر سازد.

اگرچه ساخت یک رایانهٔ کوانتومی «بحرانی از منظر رمزنگاری» \lr{(Cryptographically Relevant Quantum Computer -- CRQC)} همچنان با چالش‌های مهندسی جدی روبه‌روست، افق زمانی برآوردشده توسط بسیاری از نهادهای امنیتی و استانداردسازی (از جمله \lr{NIST} و \lr{NSA}\LTRfootnote{National Security Agency (United States)}) برای ظهور چنین سامانه‌ای، در بازهٔ $10$ تا $20$ سال آینده است. با این حال، تهدید واقعی بسیار زودتر آغاز می‌شود؛ زیرا در الگوی حملهٔ \textbf{«برداشت اکنون، رمزگشایی بعداً»} \lr{(Harvest Now, Decrypt Later -- HNDL)}، مهاجم می‌تواند داده‌های رمزشدهٔ حساس (نظیر تراکنش‌های بانکی، مکاتبات دولتی و اطلاعات هویتی) را همین امروز رهگیری و ذخیره کند تا در آیندهٔ نه‌چندان دور، پس از دسترسی به رایانهٔ کوانتومی، آن‌ها را رمزگشایی نماید. این واقعیت، ضرورت اقدام پیشدستانه را برای هر داده‌ای که ارزش محرمانگی آن بیش از افق زمانی ظهور \lr{CRQC} است -- نظیر اسناد مالی، نظامی و هویتی -- دوچندان می‌کند.

\section{چرا اکنون؟}

سه روند هم‌زمان، این پروژه را از یک تمرین آکادمیک به یک اولویت راهبردی تبدیل می‌کند:

\begin{enumerate}
\item \textbf{نهایی‌شدن استانداردهای جهانی:} در تاریخ \lr{13} آگوست \lr{2024}، مؤسسهٔ ملی استاندارد و فناوری آمریکا \lr{(NIST)} پس از فرآیندی هشت‌ساله، سه استاندارد نهایی رمزنگاری پساکوانتومی -- \lr{ML-KEM (FIPS 203)}، \lr{ML-DSA (FIPS 204)} و \lr{SLH-DSA (FIPS 205)} -- را منتشر کرد \Cite{nist2024release}، که مسیر مهاجرت جهانی را از حالت نظری به عملیاتی تبدیل کرده است.
\item \textbf{بلوغ صنعتی \lr{QKD}:} فناوری \lr{QKD} از آزمایشگاه به مقیاس ملی و بین‌قاره‌ای رسیده است؛ نمونهٔ برجستهٔ آن کریدور $2000$ کیلومتری پکن--شانگهای در چین \Cite{mdpi2023qkdn} و پیوند اخیر $303$ کیلومتری در سوئد است \Cite{arxiv2606sweden}.
\item \textbf{ورود جدی بخش بانکی:} بانک‌های پیشرو جهان نظیر \lr{HSBC} و \lr{JPMorgan Chase} دیگر در فاز آزمایشگاهی نیستند و مستقیماً معاملات واقعی و شبیه‌سازی‌شده را روی کانال‌های امن کوانتومی اجرا می‌کنند \Cite{techinformed2025banks}.
\end{enumerate}

\section{اهداف طرح}

اهداف اصلی این پروپوزال به شرح زیر است:

\begin{itemize}
\item ارائهٔ تصویری جامع و فنی از فناوری \lr{QKD} و \lr{PQC} برای تصمیم‌گیران سیاست‌گذاری و فنی کشور؛
\item طراحی معماری یک شبکهٔ آزمایشی \lr{(Pilot)} کلان‌شهری در تهران با تمرکز بر اتصال بانک مرکزی، مراکز دادهٔ بانکی و نهادهای حساس دولتی؛
\item ارائهٔ برآورد هزینهٔ تجهیزات، زیرساخت فیبر و نگهداشت بر مبنای قیمت‌های واقعی بازار جهانی؛
\item تدوین نقشهٔ راه اجرایی چندمرحله‌ای همراه با تحلیل ریسک امنیتی و پیشنهادهای کاهش آسیب‌پذیری؛
\item زمینه‌سازی برای انتقال دانش فنی، بومی‌سازی تدریجی تجهیزات و تربیت نیروی متخصص داخلی.
\end{itemize}

%=====================================================================
